% Einheit 5 von 6 – Transformationen (bank/funktionsklassen-und-eigenschaften/e5.jsonl, zusammenbau v0.1)
%% TODO Zeitmarke Einheit 5: Marken-Zeile nicht lesbar
%% TODO Zweigzeile Teil 1: was hier gelernt wird (Satz in Schülersprache aus dem Einheitstitel) – Einheit 5
\einheitenkopf[e5]{Einheit 5 von 6 · Transformationen}
\zweigzeile{Abitur GK}
\setcounter{aufgabe}{18}

%% TODO Titel: Ich-kann-Satz fehlt in der Bank (Platzhalter: Kettenname) – Nr. 19
%% TODO Anweisung: ein Satz über der ersten Teilaufgabe fehlt in der Bank – Nr. 19
\begin{aufgabe}{Transformationen}
%% TODO \rechenplatz{n}: Schrittzahl des Grundfalls fehlt in der Bank (nur bei fester Schreibform, 2.3 b)
\begin{teile}
\teil $f(x)$ und $f(x - 4)$ – innen oder außen? \\ \kreuz{innen: x-Richtung} \\ \kreuz{außen: y-Richtung}
\teil $g(x) = f(x) + 3$ – wie entsteht der Graph von $g$ aus dem von $f$?
\teil $f(x) = x^2$ und $g(x) = (x + 2)^2$ – wie geht der Graph von $g$ aus dem von $f$ hervor?
\teil Der Graph von $f(x) = x^2$ ist gezeichnet. Skizziere den Graphen von $g(x) = (x - 1)^2 - 2$.

\begin{ksys}[xmin=-4,xmax=4,ymin=-3,ymax=5]\parabel{1}{0}{0}{f}\end{ksys}
\teil $f(x) = x^3 - 12x$ hat den Hochpunkt $H(-2 | 16)$. Der Graph wird um $3$ nach rechts und um $5$ nach oben verschoben. Term und Hochpunkt des Bildes?
\teil Der Graph von $g$ entsteht aus dem von $f$ durch Streckung in y-Richtung; aus $P(2 | 3)$ wird $P'(2 | 7{,}5)$. Streckfaktor?
\teil $f(x) = x^3$ und $g(x) = f(x) + 3$. Aussage: Beide Graphen haben an der Stelle $2$ dieselbe Steigung. Stimmt das?
\teil $f(x) = 3\,\mathrm{sin}(2x) + 1$ – Amplitude, Periode, Mittellinie?
\teil Die obere Randlinie eines Schildes ist $f(x) = -0{,}1x^2 + 2$ für $-4 \le x \le 4$ (1 Längeneinheit entspricht 1 dm). Die untere Randlinie entsteht durch Spiegelung an der x-Achse. Term der unteren Randlinie und Höhe des Schildes in der Mitte?
\teil $g(x) = f(x - 2) + 1$ und $f'(3) = 4$. Wie groß ist $g'(5)$?
\teil $f$ ist in $\mathbb{R}$ definiert und hat die Wertemenge $[-1; 4[$; $h(x) = -3 \cdot f(x + 2) + 2$. Gib die Wertemenge von $h$ an und begründe \hfill (Abitur 2026 LK).
\end{teile}
\end{aufgabe}

%% TODO Titel: Ich-kann-Satz fehlt in der Bank (Platzhalter: Kettenname) – Nr. 20
%% TODO Anweisung: ein Satz über der ersten Teilaufgabe fehlt in der Bank – Nr. 20
\begin{aufgabe}{Transformationen – Fehler finden, Begründen, Darstellungswechsel, Anwendung}
\begin{teile}
\teil Ich finde den Fehler: „Der Graph von $g(x) = f(x + 3)$ ist der um $3$ nach rechts verschobene Graph von $f$.“
\teil Begründe: Der Graph von $g(x) = f(x - 2)$ ist um $2$ nach rechts verschoben, obwohl im Term ein Minus steht.
\teil Die Graphen von $f(x) = x^2$ und $g$ sind gezeichnet. Lies ab: Wie entsteht $g$ aus $f$? Gib den Term von $g$ an.

\begin{ksys}[xmin=-3,xmax=5,ymin=-2,ymax=5,ablesen]\parabel{1}{0}{0}{f}\parabel{1}{2}{-1}{g}\end{ksys}
\teil Wasserstand in einem Hafen: $h(t) = 2\,\mathrm{cos}(0{,}5t) + 5$ ($h$ in m, $t$ in Stunden). Was bedeutet die Periode? Höchster und niedrigster Wasserstand?
\end{teile}
\end{aufgabe}
